% ===== PRÉAMBULE CANONIQUE QCM — MATHÉMATIQUES =====
% Sert à compiler controle.tex ET à la revalidation automatique du JSON
% (valider_qcm.py). NE PAS MODIFIER : packages figés.
\documentclass[11pt,a4paper]{article}
\usepackage[utf8]{inputenc}\usepackage[T1]{fontenc}\usepackage{lmodern}
\usepackage[french]{babel}
\usepackage{amsmath,amssymb,mathtools}\usepackage{icomma}
\usepackage{siunitx}\sisetup{locale=FR}
\usepackage{xcolor}\usepackage{colortbl}
\usepackage{tikz}\usepackage{pgfplots}\pgfplotsset{compat=1.18}
\usepackage{enumitem}\usepackage{array,booktabs}
\usepackage[a4paper,margin=2.2cm]{geometry}
\usetikzlibrary{arrows.meta,calc,angles,quotes,positioning,patterns,backgrounds,shapes.geometric,intersections,babel}
\newcommand{\R}{\mathbb{R}}\newcommand{\N}{\mathbb{N}}\newcommand{\Z}{\mathbb{Z}}
\newcommand{\Q}{\mathbb{Q}}\newcommand{\C}{\mathbb{C}}\newcommand{\ds}{\displaystyle}
\newcommand{\vv}[1]{\vec{#1}}
% ===================================================
\begin{document}

\begin{center}
{\large\textbf{QCM concours --- Mathématiques --- Fiche 6 : Les limites}}\par
\textit{10 questions, niveau concours. Une seule proposition correcte par question.}
\end{center}
\vspace{6pt}\hrule\vspace{10pt}

% ------------------------------------------------------------------
\noindent\textbf{MATH-F06-Q01 --- Limite en un point : forme $\frac00$ par factorisation}\par
Que vaut $\ds\lim_{x\to 3}\frac{x^{2}-9}{x^{3}-27}$ ?
\begin{enumerate}[label=\Alph*)]
\item $\dfrac{1}{3}$
\item $\dfrac{2}{9}$
\item $\dfrac{2}{3}$
\item $\dfrac{1}{9}$
\end{enumerate}
\textit{Bonne réponse : B}\par
La forme directe est $\frac00$. On factorise avec $a^{2}-b^{2}$ et $a^{3}-b^{3}$ :
$x^{2}-9=(x-3)(x+3)$ et $x^{3}-27=(x-3)(x^{2}+3x+9)$. Après simplification de $(x-3)$,
\[ \lim_{x\to 3}\frac{x^{2}-9}{x^{3}-27}=\lim_{x\to 3}\frac{x+3}{x^{2}+3x+9}=\frac{6}{9+9+9}=\frac{6}{27}=\frac{2}{9}. \]
\textit{Pièges :}
\begin{itemize}
\item A : applique la règle du terme dominant ($x^{2}/x^{3}=1/x\to 1/3$) comme à l'infini, alors qu'on est en un point fini.
\item C : factorise le dénominateur avec un mauvais signe ($x^{2}-3x+9$), d'où $6/(9-9+9)=2/3$.
\item D : simplifie $6/27$ en $1/9$ (dénominateur doublé par erreur).
\end{itemize}
\vspace{8pt}\hrule\vspace{8pt}

% ------------------------------------------------------------------
\noindent\textbf{MATH-F06-Q02 --- Asymptote verticale : règle $a/0^{\pm}$ (limite latérale)}\par
Que vaut $\ds\lim_{x\to 1^{+}}\frac{x^{2}-3x}{x^{2}-1}$ ?
\begin{enumerate}[label=\Alph*)]
\item $-\infty$
\item $+\infty$
\item $\dfrac{1}{2}$
\item $1$
\end{enumerate}
\textit{Bonne réponse : A}\par
Le numérateur tend vers $1-3=-2\neq 0$. Le dénominateur $x^{2}-1=(x-1)(x+1)$ s'annule en $1$ :
quand $x\to 1^{+}$, $(x-1)\to 0^{+}$ et $(x+1)\to 2$, donc $x^{2}-1\to 0^{+}$. Par la règle des signes,
\[ \lim_{x\to 1^{+}}\frac{x^{2}-3x}{x^{2}-1}=\frac{-2}{0^{+}}=-\infty. \]
\textit{Pièges :}
\begin{itemize}
\item B : confond le côté (à gauche $x^{2}-1\to 0^{-}$ donnerait $+\infty$) ou oublie que le numérateur est négatif.
\item C : « simplifie » abusivement $\dfrac{x(x-3)}{(x-1)(x+1)}$ en $\dfrac{x}{x+1}$ puis remplace $x=1$.
\item D : confond avec la limite à l'infini (degrés égaux $x^{2}/x^{2}\to 1$).
\end{itemize}
\vspace{8pt}\hrule\vspace{8pt}

% ------------------------------------------------------------------
\noindent\textbf{MATH-F06-Q03 --- Asymptote verticale : limite bilatérale inexistante}\par
Que vaut $\ds\lim_{x\to 2}\frac{x^{2}-3x}{x^{2}-4}$ ?
\begin{enumerate}[label=\Alph*)]
\item $+\infty$
\item $-\infty$
\item La limite n'existe pas
\item $1$
\end{enumerate}
\textit{Bonne réponse : C}\par
Le numérateur tend vers $4-6=-2\neq 0$ ; le dénominateur $(x-2)(x+2)$ s'annule en $2$. Il faut étudier
les deux côtés :
\[ \lim_{x\to 2^{-}}\frac{x^{2}-3x}{x^{2}-4}=\frac{-2}{0^{-}}=+\infty,\qquad
   \lim_{x\to 2^{+}}\frac{x^{2}-3x}{x^{2}-4}=\frac{-2}{0^{+}}=-\infty. \]
Les deux limites latérales sont opposées : la limite (bilatérale) \textbf{n'existe pas}.
\textit{Pièges :}
\begin{itemize}
\item A : ne considère que la limite à gauche ($x\to 2^{-}$ : $-2/0^{-}=+\infty$).
\item B : ne considère que la limite à droite ($x\to 2^{+}$ : $-2/0^{+}=-\infty$).
\item D : calcule la limite à l'infini (degrés égaux $x^{2}/x^{2}\to 1$) au lieu de la limite en $2$.
\end{itemize}
\vspace{8pt}\hrule\vspace{8pt}

% ------------------------------------------------------------------
\noindent\textbf{MATH-F06-Q04 --- Fonction rationnelle à l'infini : degrés égaux}\par
Que vaut $\ds\lim_{x\to +\infty}\frac{9x^{4}-2x^{2}+x}{-6x^{4}+5x-3}$ ?
\begin{enumerate}[label=\Alph*)]
\item $\dfrac{3}{2}$
\item $-\dfrac{2}{3}$
\item $-2$
\item $-\dfrac{3}{2}$
\end{enumerate}
\textit{Bonne réponse : D}\par
À l'infini, la limite d'une fonction rationnelle est celle du rapport des termes de plus haut degré :
\[ \lim_{x\to +\infty}\frac{9x^{4}-2x^{2}+x}{-6x^{4}+5x-3}=\lim_{x\to +\infty}\frac{9x^{4}}{-6x^{4}}=\frac{9}{-6}=-\frac{3}{2}. \]
\textit{Pièges :}
\begin{itemize}
\item A : oublie le signe $-$ du terme dominant du dénominateur ($-6x^{4}$).
\item B : inverse le rapport ($b_m/a_n=-6/9=-2/3$ au lieu de $a_n/b_m$).
\item C : « somme tous les coefficients » comme si l'on remplaçait $x=1$ : $(9-2+1)/(-6+5-3)=8/(-4)=-2$.
\end{itemize}
\vspace{8pt}\hrule\vspace{8pt}

% ------------------------------------------------------------------
\noindent\textbf{MATH-F06-Q05 --- Fonction rationnelle à l'infini : degrés différents}\par
Que vaut $\ds\lim_{x\to -\infty}\frac{-x^{4}+2x}{x^{3}-5}$ ?
\begin{enumerate}[label=\Alph*)]
\item $-\infty$
\item $+\infty$
\item $-1$
\item $0$
\end{enumerate}
\textit{Bonne réponse : B}\par
On garde les termes dominants : $\dfrac{-x^{4}}{x^{3}}=-x$. Or, quand $x\to -\infty$, $-x\to +\infty$. Donc
\[ \lim_{x\to -\infty}\frac{-x^{4}+2x}{x^{3}-5}=\lim_{x\to -\infty}(-x)=+\infty. \]
\textit{Pièges :}
\begin{itemize}
\item A : erreur de double signe --- écrit $-x$ puis conclut $-\infty$, en oubliant que $-(-\infty)=+\infty$.
\item C : croit les degrés égaux et prend le rapport des coefficients dominants $-1/1=-1$.
\item D : croit le degré du numérateur inférieur à celui du dénominateur (limite nulle).
\end{itemize}
\vspace{8pt}\hrule\vspace{8pt}

% ------------------------------------------------------------------
\noindent\textbf{MATH-F06-Q06 --- Asymptotes horizontale et verticale : point d'intersection}\par
La courbe représentative de $f(x)=\dfrac{2x+4}{x-3}$ admet une asymptote verticale et une asymptote
horizontale. Quelles sont les coordonnées de leur point d'intersection ?
\begin{enumerate}[label=\Alph*)]
\item $(3\,;\,2)$
\item $(-3\,;\,2)$
\item $(2\,;\,3)$
\item $\left(3\,;\,\dfrac{1}{2}\right)$
\end{enumerate}
\textit{Bonne réponse : A}\par
\textbf{Asymptote verticale} : le dénominateur s'annule en $x=3$ (le numérateur y vaut $10\neq 0$), donc
$x=3$. \textbf{Asymptote horizontale} : $\ds\lim_{x\to\pm\infty}\frac{2x+4}{x-3}=\frac{2x}{x}=2$, donc $y=2$.
Le point d'intersection des deux droites $x=3$ et $y=2$ est $(3\,;\,2)$.
\textit{Pièges :}
\begin{itemize}
\item B : résout $x-3=0$ en $x=-3$ (erreur de signe sur l'asymptote verticale).
\item C : intervertit l'abscisse et l'ordonnée du point.
\item D : prend l'asymptote horizontale comme le rapport dominant inversé $\tfrac12$ au lieu de $2$.
\end{itemize}
\vspace{8pt}\hrule\vspace{8pt}

% ------------------------------------------------------------------
\noindent\textbf{MATH-F06-Q07 --- Limite à l'infini d'une somme de termes}\par
Que vaut $\ds\lim_{x\to +\infty}\left(\frac{6x^{2}+1}{2x^{2}-x}-\frac{4}{x}+1\right)$ ?
\begin{enumerate}[label=\Alph*)]
\item $3$
\item $\dfrac{4}{3}$
\item $0$
\item $4$
\end{enumerate}
\textit{Bonne réponse : D}\par
On calcule la limite de chaque terme séparément :
\[ \frac{6x^{2}+1}{2x^{2}-x}\to\frac{6x^{2}}{2x^{2}}=3,\qquad \frac{4}{x}\to 0,\qquad 1\to 1. \]
D'où $\ds\lim_{x\to +\infty}\left(\frac{6x^{2}+1}{2x^{2}-x}-\frac{4}{x}+1\right)=3-0+1=4.$
\textit{Pièges :}
\begin{itemize}
\item A : oublie le terme $+1$ et ne garde que la limite de la fraction ($=3$).
\item B : inverse le rapport dominant ($\tfrac{2}{6}=\tfrac13$), d'où $\tfrac13-0+1=\tfrac43$.
\item C : croit que $\tfrac{4}{x}\to 4$ au lieu de $0$, d'où $3-4+1=0$.
\end{itemize}
\vspace{8pt}\hrule\vspace{8pt}

% ------------------------------------------------------------------
\noindent\textbf{MATH-F06-Q08 --- Radical en un point : forme $\frac00$ (quantité conjuguée)}\par
Que vaut $\ds\lim_{x\to 1}\frac{\sqrt{x+3}-2}{x-1}$ ?
\begin{enumerate}[label=\Alph*)]
\item $\dfrac{1}{2}$
\item $0$
\item $\dfrac{1}{4}$
\item $\dfrac{1}{8}$
\end{enumerate}
\textit{Bonne réponse : C}\par
Forme $\frac00$. On multiplie haut et bas par le conjugué $\sqrt{x+3}+2$ :
\[ \frac{\sqrt{x+3}-2}{x-1}=\frac{(x+3)-4}{(x-1)\bigl(\sqrt{x+3}+2\bigr)}=\frac{x-1}{(x-1)\bigl(\sqrt{x+3}+2\bigr)}=\frac{1}{\sqrt{x+3}+2}. \]
Le facteur $(x-1)$ se simplifie, donc
$\ds\lim_{x\to 1}\frac{\sqrt{x+3}-2}{x-1}=\frac{1}{\sqrt{4}+2}=\frac{1}{4}.$
\textit{Pièges :}
\begin{itemize}
\item A : par l'Hospital, dérive $\sqrt{x+3}$ en $\dfrac{1}{\sqrt{x+3}}$ (oubli du $\frac12$), d'où $1/\sqrt{4}=1/2$.
\item B : conclut $0$ pour la forme $\frac00$ sans rationaliser.
\item D : évalue $\dfrac{1}{2\sqrt{x+3}}$ en remplaçant $\sqrt{x+3}$ par $x+3=4$, soit $1/8$.
\end{itemize}
\vspace{8pt}\hrule\vspace{8pt}

% ------------------------------------------------------------------
\noindent\textbf{MATH-F06-Q09 --- Forme $\frac00$ : factorisation d'une somme de cubes}\par
Que vaut $\ds\lim_{x\to -2}\frac{x^{3}+8}{x^{2}-4}$ ?
\begin{enumerate}[label=\Alph*)]
\item $-3$
\item $-1$
\item $-2$
\item $3$
\end{enumerate}
\textit{Bonne réponse : A}\par
Forme $\frac00$. On factorise avec $a^{3}+b^{3}$ et $a^{2}-b^{2}$ :
$x^{3}+8=(x+2)(x^{2}-2x+4)$ et $x^{2}-4=(x-2)(x+2)$. Après simplification de $(x+2)$,
\[ \lim_{x\to -2}\frac{x^{3}+8}{x^{2}-4}=\lim_{x\to -2}\frac{x^{2}-2x+4}{x-2}=\frac{4+4+4}{-4}=\frac{12}{-4}=-3. \]
\textit{Pièges :}
\begin{itemize}
\item B : mauvais signe dans le facteur ($x^{2}+2x+4$ au lieu de $x^{2}-2x+4$), d'où $4/(-4)=-1$.
\item C : applique la règle du terme dominant $x^{3}/x^{2}=x\to -2$ au lieu de factoriser.
\item D : erreur de signe final ($12/4=3$ en oubliant que le dénominateur vaut $-4$).
\end{itemize}
\vspace{8pt}\hrule\vspace{8pt}

% ------------------------------------------------------------------
\noindent\textbf{MATH-F06-Q10 --- Limite trigonométrique remarquable}\par
Que vaut $\ds\lim_{x\to 0}\frac{\sin(7x)}{2x}$ ?
\begin{enumerate}[label=\Alph*)]
\item $\dfrac{1}{2}$
\item $\dfrac{7}{2}$
\item $\dfrac{2}{7}$
\item $7$
\end{enumerate}
\textit{Bonne réponse : B}\par
On fait apparaître la limite de référence $\dfrac{\sin u}{u}\to 1$ en multipliant et divisant par $7x$ :
\[ \frac{\sin(7x)}{2x}=\underbrace{\frac{\sin(7x)}{7x}}_{\to\,1}\times\frac{7x}{2x}=\frac{\sin(7x)}{7x}\times\frac{7}{2}\ \xrightarrow[x\to 0]{}\ 1\times\frac{7}{2}=\frac{7}{2}. \]
\textit{Pièges :}
\begin{itemize}
\item A : approxime $\sin(7x)\approx x$ en oubliant le facteur $7$ interne, d'où $x/(2x)=1/2$.
\item C : inverse le rapport ($2/7$ au lieu de $7/2$).
\item D : oublie le $2$ du dénominateur ($\sin(7x)\approx 7x$, puis $7x/x=7$).
\end{itemize}
\vspace{8pt}\hrule\vspace{8pt}

\end{document}
